\section{The exact shifted-zeta evolution}\label{sec:shift}

\subsection{The transfer and its causal generator}

Define
\begin{equation}
 \xi(s)=\tfrac12s(s-1)\pi^{-s/2}\Gamma(s/2)\zeta(s),
 \qquad H(p)=\xi(\tfrac12+p).
\end{equation}
For $0<\omega\leq1/2$, set
\begin{equation}\label{eq:transfer}
 K_\omega(p)=\frac{H(p-\omega)}{H(p+\omega)},\qquad
 a_\omega(p)=\frac{H'(p-\omega)}{H(p-\omega)}
             +\frac{H'(p+\omega)}{H(p+\omega)}.
\end{equation}
Logarithmic differentiation gives
$\partial_\omega K_\omega=-a_\omega K_\omega$.
The functional equation gives $|K_\omega(\ii\tau)|=1$ wherever
the ratio is defined. This boundary identity alone does not give
a bounded causal transfer on the unweighted whole line.

Instead take inverse Laplace transforms on $\Ree p=\sigma>1$,
where the denominator is nonzero and the Dirichlet series converge.
If $k_\omega$ is the causal kernel, define
\begin{equation}
 V_{\omega,L}f=(k_\omega*F)|_{I_L},\qquad V_{0,L}=I.
\end{equation}
For positive shift the kernel is locally integrable and the
finite-window operator is bounded. The exact Beta-kernel, integer
comb and rational correction are given in
\cite[Section~S1.4]{Supplement}; positivity of the Beta/comb part
does not imply contraction of the full transfer.

Let $G_{\omega,L}$ be the compressed causal realization of
$a_\omega$ on the same Laplace line. Its symmetric quadratic form is
\begin{equation}\label{eq:Qomega}
 Q_{\omega,L}[g]=\Ree\ip g{G_{\omega,L}g}.
\end{equation}
This notation is interpreted as a form when the operator pairing
is not separately defined. The natural form domain has logarithmic
Fourier weight as in Section~\ref{sec:arithmetic}; it need not be
the full operator domain of the nonsymmetric generator.
Smooth compactly supported tests and their positive-shift
evolutions suffice for the identities below.

\subsection{The equation to be realized}

Causality permits the scalar transfer identity to be compressed to
the finite interval: an intermediate value beyond its right endpoint
cannot return to the interval, and the input has no values before
its left endpoint. The exact equation is
\begin{equation}\label{eq:prescribed-flow}
 \boxed{\partial_\omega V_{\omega,L}
       =-G_{\omega,L}V_{\omega,L},\qquad V_{0,L}=I.}
\end{equation}
It is the boxed equation adjacent to the energy derivative in the
shared background~\cite{SharedBackground}. It holds on admissible
inputs, with a strong initial limit and a right derivative at zero
on smooth tests. No operator-norm derivative at zero is asserted.

The zero-shift symbol shows the matching problem explicitly:
\begin{align}\label{eq:generator}
 a_0(p)={}&w_0+
 2\int_0^\infty\nGamma(u)(1-e^{-pu})\dd u
 +\frac2{p+1/2}+\frac2{p-1/2}\notag\\
 &-2\sum_{n\geq2}\frac{\Lambda(n)}{\sqrt n}e^{-p\log n}.
\end{align}
Its real quadratic form is exactly $Q_{0,L}$, including the fixed
local constant, the pole response and the prime delays.
At positive shift the prime term in $a_\omega$ is
\[
 -2\sum_{n\geq2}\frac{\Lambda(n)}{\sqrt n}
       \cosh(\omega\log n)e^{-p\log n}.
\]
The transfer itself contains the fuller multiplicative integer comb,
not just the prime-power part of its logarithmic derivative.

Thus \eqref{eq:prescribed-flow} supplies a detailed prescription.
A physical construction must reproduce every component of
\eqref{eq:generator}, and its finite-shift continuation, with these
coefficients. Choosing a contour with a suggestive gamma spectrum
addresses only one component.

\subsection{The accumulated defect and the central limit}

Put
\begin{equation}\label{eq:defect}
 D_{\omega,L}=I-V_{\omega,L}^*V_{\omega,L}.
\end{equation}
The product rule at positive shift yields the second evolution
equation
\begin{equation}\label{eq:defect-flow}
 \boxed{\partial_\omega D_{\omega,L}
 =V_{\omega,L}^*(G_{\omega,L}^*+G_{\omega,L})V_{\omega,L},
 \qquad D_{0,L}=0.}
\end{equation}
The right side is the pullback of the symmetric form $2Q_{\omega,L}$.
Equivalently one uses the bounded positive-shift product-rule
expression
$(G_{\omega,L}V_{\omega,L})^*V_{\omega,L}
+V_{\omega,L}^*(G_{\omega,L}V_{\omega,L})$.
This avoids assuming that the separate adjoint generator acts
on every evolved input.

For $f\in C_c^\infty(I_L)$, integration and the initial limit give
\begin{align}
 \ip f{D_{\omega,L}f}
 &=\norm f^2-\norm{V_{\omega,L}f}^2
 =2\int_0^\omega Q_{s,L}[V_{s,L}f]\dd s,
 \label{eq:cumulative}\\
 Q_{0,L}[f]
 &=\lim_{\omega\downarrow0}
        \frac{\ip f{D_{\omega,L}f}}{2\omega}.
 \label{eq:central-limit}
\end{align}
The running shift $s$ changes both the form and the input to it.
The analytic justification in the shared background uses the
gamma-ratio bounds
$|\partial_\omega^jK_\omega(\sigma+\ii\tau)|
\leq C_\sigma(1+\log(2+|\tau|))^j(2+|\tau|)^{-\omega}$
for $j=0,1,2$: these give positive-shift smoothing and
differentiability, and dominated convergence on smooth inputs at
zero. No inverse of $V_{\omega,L}$ is needed.

The positivity target is $D_{\omega,L}\succeq0$, or equivalently
$\norm{V_{\omega,L}}\leq1$. Requiring all the instantaneous
$Q_{s,L}$ to be positive would suffice, but is a stronger demand
than positivity of their accumulated pullback in
\eqref{eq:cumulative}. This distinction determines the proposed
role of the reflected network in Section~\ref{sec:storage}.
